% Status: Schlussfolgerung aus der Analyse (4b), ausgeschrieben.
\chapter{Schlussfolgerung aus der Analyse}
\label{ch:schluss}

\section{Synthese der drei Hypothesen}
Die drei Hypothesen ergeben ein konsistentes, aber bewusst begrenztes Bild. H1
zeigt eine bessere Prognosequalität von Polymarket nur in einem klar definierten
Vergleichsscope: In der nationalen Tagesreihe trägt Polymarket den tieferen
mittleren Brier-Verlust (0.230 gegen 0.332 von FiveThirtyEight), und auf
Bundesstaaten-Ebene stützen 262 von 285 Zeilen sowie 9 von 9 Staaten den Vorteil
(exakter Binomialtest, p = 0.0020). Die breite Überlegenheit bleibt jedoch
ungedeckt, und das volle State-Date-Panel ist ein Gegenbeispiel. H2 zeigt plausible
Tagesreaktionen auf die sieben vorab kuratierten Ereignisse, am deutlichsten beim
Attentat auf Trump (+7.2 Prozentpunkte), während die neutral kuratierten Ereignisse
erwartungsgemäss streuen. H3 findet das klarste prädiktive Timing-Muster im
obersten Wallet-Tier (Lag-1-Korrelation 0.186, minimaler Granger-p-Wert 0.0012 über
1216 ausgerichtete Zeilen). Zusammengenommen stützen die Befunde begrenzte,
scope-spezifische Diagnosen informationeller Effizienz, keine universelle Aussage.

\section{Stützende Evidenz ausserhalb der Hypothesen}
Zwei weitere Beobachtungen fügen sich in dieses Bild. Erstens deckt sich der
paper-only Arbitrage-Befund (Kapitel~\ref{ch:loesung}) mit weitgehend effizienten
Märkten: Saubere strukturelle Arbitragen sind selten, die wenigen echten
Cross-Venue-Fälle sind kleine Carry-Trades, und die kaum robust ausnutzbare
Ineffizienz ist selbst ein Effizienz-Indiz. Zweitens lag Polymarket in der
Schweizer Referendums-Fallstudie als binäre Wahrscheinlichkeit auf der richtigen
Seite des Resultats; dies ist jedoch ein transparenter Proxy und ein einzelner,
frischer Datenpunkt ausserhalb der US-Wahl, keine Verallgemeinerung.

\section{Einordnung in die Literatur}
Der begrenzte H1-Vorteil ist mit \citet{zotero_poly_002} vereinbar, die einen
Polymarket-Vorteil besonders in Swing States berichten, zugleich aber zur Vorsicht
bei Verallgemeinerung mahnen. Das Gegenbeispiel im vollen Panel passt zur kritischen
Lesart von \citet{zotero_poly_009}, wonach Prognosemärkte reflexiv statt eine
``Truth Machine'' sind, sowie zur gemischten plattformübergreifenden Evidenz für
2024 \citep{lit_poly24_accuracy}. Die Vorsicht bei H3 deckt sich mit
\citet{zotero_poly_001}, deren Grosskonto-Befund heterogene Erwartungen statt
Manipulation nahelegt, und mit der Insider-Rahmung von \citet{lit_kyle_001} und
\citet{zotero_delvecchio_001}. Insgesamt verarbeiten die Märkte Information
innerhalb von Grenzen gut, vereinbar mit einer halbstrengen Effizienzlesart
\citep{lit_emh_001}, ohne dass Preis und Wahrscheinlichkeit unbesehen
gleichzusetzen wären \citep{lit_manski_2006}.

\section{Folgerung für die Forschungsfrage}
Aus der Analyse folgt, dass Polymarket im untersuchten Kontext Information
nachweisbar gut verarbeitet, aber nur in klar abgegrenzten Fällen messbar besser
als traditionelle Quellen. Diese begrenzte, doch belastbare Diagnose ist der
eigentliche Beitrag: Sie trennt Gestütztes sauber von Nicht-Gestütztem. Genau das
motiviert den folgenden Lösungsteil --- ein Werkzeug, das solche begrenzten,
deterministisch geprüften Hinweise sichtbar und für Menschen prüfbar macht.
